\chapter{Figures}\label{ch:figures}
\begin{goals}
\textbf{In this chapter} you turn a net into a page for a report: the net as vectors, a legend, the
statistics and a caption, laid out as you like and exported as PDF or PNG.
\end{goals}

A net on the desk is drawn for working: it follows the theme, shows the selection, and changes as you
change things. A \term{figure} is a page made \emph{from} a net, for someone else to read. wuffle builds it
with compo, the GCU page compositor, and keeps it in the project like a dataset or a diagram.

\section{Making one}\label{sec:new-figure}

Make the net you want active, then choose \menu{Figure>New figure from the active net}. A tab called
\ui{<net> · figure} opens with an A4 landscape page holding five frames:

\begin{description}
\item[The title] the net's name.
\item[The net] drawn as vectors, not as a picture, so it stays sharp at any size and in the PDF. It is
drawn on a white page whatever the desk's theme (the layers keep their colours, the graticule and
arrows are ink), and it shows the data, not the selection.
\item[The legend] each visible layer with its colour and its count, \ui{tocher (200)}.
\item[The statistics] a table: for each layer, n, the axial Fisher mean, $\kappa$, $\alpha_{95}$ and
Woodcock's K (\cref{ch:stats}).
\item[The caption] what the reader is looking at, in words: the projection and hemisphere, the rotation if
any (\cref{ch:rotation}), how the contours were made, the notation (\enquote{Planes as dip direction/dip,
lines as trend→plunge}), the lineage of any derived dataset on the net, and \enquote{Made with wuffle} with
the date.
\end{description}

Only nets make figures. The menu item is greyed out while a rose diagram or a fabric plot is active.

\shot{figure-editor}{A figure from the contoured Tocher net. The tree lists it under \ui{figures}; the
frames panel on the right lists the page's five frames.}

The caption is generated, so it can't disagree with the net. For the Tocher net contoured by Kamb's
method it reads:

\begin{quote}\small
Equal-area (Schmidt), lower hemisphere · Kamb contours at 2σ intervals (k = 3). Planes as dip
direction/dip, lines as trend→plunge. Made with wuffle, 26/11/2026.
\end{quote}

\section{Laying out the page}

The figure tab is compo's editor. Click a frame to edit it, drag it to move it (it snaps to edges, centres
and equal gaps), and pull its handles to resize it; the arrow keys nudge. A right-click on a frame offers
duplicate, group (when two or more frames are selected), the stacking order, lock and delete. The bar above the page sets the paper size,
orientation, zoom and grid, and adds new frames, such as text.

Inside the figure tab \key{Ctrl}+\key{Z} is compo's own undo, step by step through your edits to the
page. Outside it, the figure's changes are part of the project's history like any other (\cref{sec:undo}).

\section{Keeping it in step with the net}\label{sec:figure-refresh}

A figure is a copy of the net at the moment you made it. When the net changes (more data, a different
contour method, a rotation), click \ui{↻ from the net} in the figure's bar. The five frames are made again
from the diagram as it is now; where you put them, how big you made them and any frames you added are
kept.

If the diagram the figure was made from has been removed, \ui{↻ from the net} says so and changes
nothing.

\section{Export}\label{sec:figure-export}

\begin{description}
\item[\ui{Export PDF}] the page as a vector PDF, \file{<name>.pdf}. The net is vectors in it too.
\item[\ui{Export PNG}] the page as a picture at 200 dpi, \file{<name>.png}.
\end{description}

In the file name, each run of characters other than letters, digits, underscores, dots and hyphens becomes one \code{\_}: the figure
\ui{Net 1 · figure} exports as \file{Net\_1\_figure.pdf}.

\section{Figures in the project}\label{sec:figures-project}

Figures are saved with the project (\cref{sec:project-file}). The tree lists them under \ui{figures}: click
one to open its tab, or right-click to rename or delete it. The \menu{Figure} menu lists them too.

In the console each figure is a file under \file{/project/figures} (\cref{ch:project-fs}):

\begin{description}
\item[\file{<name>.compo.json}] the page itself, as compo's JSON. Writing it replaces the page; a new name
is a new figure; \code{mv} renames it and \code{rm} deletes it.
\item[\file{<name>.pdf}] the page as a PDF, made when you read it. \code{cp} it to \file{/home} to keep a
copy in the project.
\end{description}

\begin{tryit}
\begin{enumerate}
\item Open the Tocher sample and, in the tree, tick \ui{contours} and \ui{mean + cone} on \ui{tocher}.
\item Choose \menu{Figure>New figure from the active net}. The figure opens in its own tab.
\item Drag the legend frame down the page.
\item In the console, type \code{contour fisher}. The net on the desk changes; the figure doesn't.
\item Back in the figure, click \ui{↻ from the net}. The contours change and the caption now reads
\enquote{Fisher-kernel density contours}; the legend stays where you put it.
\item Click \ui{Export PDF}.
\end{enumerate}
\end{tryit}
